% SPDX-FileCopyrightText: 2026 Hasan Melih Akbulut
% SPDX-License-Identifier: CC-BY-4.0
%
% The numbers this paper's argument turns on are registered in
% paper/claims.yaml and checked by paper/check_claims.py, which runs in
% CI through .github/workflows/checks.yml -> scripts/ci_local.sh paper.
%
% Corrected 2026-09-12. This read "configured to run in CI and has not
% been able to since 2026-09-03". Both halves were wrong. 53 of the 61
% workflow runs in that window executed; what failed was one job, on a
% missing PyYAML, and one day of billing refusals was read backwards
% across a week. The claim checker has since run green on a runner, and
% the paper itself built there -- commit cf2ece6 records "build the
% paper: ok". The limitations section is still where the OBSERVED answer
% is stated, and it still says what this file cannot check locally:
% there is no TeX on the development machine, so the local build skips.
%
% The registry is NOT a census of every number printed here. It holds
% the figures the argument turns on and nothing else, so it is always
% smaller than the body's stock of numeric literals, and it grows when a
% claim is added rather than when a number is. That is the property; the
% size is not. An earlier version of this header printed both as counts
% -- "the body prints about 174 numeric literals and the registry holds
% 27" -- and the second was already stale when it was read, because a
% count of a file that grows goes stale silently while still reading as
% a fact, and nothing was checking it. Naming a count here would only
% move the same defect one line down, so this header names none.
% `python3 paper/check_claims.py --list` prints the registry as
% it stands, and the checker's own summary says how many of its claims
% are re-derived by a command and how many are read once by hand from a
% named artefact.

\documentclass[11pt,a4paper]{article}

\usepackage[T1]{fontenc}
\usepackage[utf8]{inputenc}
\usepackage{lmodern}
\usepackage[margin=1in]{geometry}
\usepackage{booktabs}
\usepackage{array}
\usepackage{amsmath}
\usepackage{siunitx}
\usepackage{microtype}
\usepackage{tabularx}
\usepackage[hidelinks]{hyperref}   % loaded last, as it patches others

% group-digits=integer, or group-minimum-digits puts a comma inside every
% four-digit DECIMAL fraction: "1.226,2 ns". detect-weight is a siunitx
% v2 key; v3 spells it text-series-to-math.
\sisetup{group-separator={,}, group-minimum-digits=4,
         group-digits=integer, text-series-to-math=true}

\newcolumntype{L}[1]{>{\raggedright\arraybackslash}p{#1}}
\newcommand{\file}[1]{\texttt{\small #1}}

\title{Triplicated in the Netlist, Adjacent in the Die:\\
Instrumenting a Redundancy Claim Across an Open 130\,nm Flow}

\author{Hasan Melih Akbulut\thanks{Independent researcher.
Correspondence: \texttt{57303760+Melihakbulut221@users.noreply.github.com}.
All sources, harnesses and records are open: see Appendix~\ref{app:artefacts}.}}

\date{9 September 2026}

\begin{document}
\maketitle

\begin{abstract}
We report two levels at which a triple-modular-redundancy claim about a
130\,nm open-PDK neuromorphic system-on-chip was false while every check
pointed at it returned green, and we name the instrument that was blind
to each.
\textbf{No silicon exists, no radiation testing has been performed, and
neither open PDK involved publishes total-ionising-dose or single-event
data}; every reliability figure here is a seeded fault-injection
measurement on RTL and on placed netlists, and the system-on-chip layout
fails 3{,}088 setup endpoints at the slow corner and closes at
\SI{25}{\nano\second}, or \SI{24}{\nano\second} after a repair we report
in full including the arm that made it worse. Its geometric decks had been switched off for
three years' worth of runs on the strength of one experiment; we ran
them, and \textbf{of \num{39969214} Magic error \emph{boxes} and
\num{9668} KLayout \emph{markers}, zero fall outside the vendor SRAM
macros}. Neither headline is a count of distinct violations and we
caveat both: Magic's own report ends by saying to divide it by 3 or 4,
and \num{3918} of the KLayout markers are one population of shapes
reported a second time by a second rule, leaving \num{5750} distinct.
Every one of those divisors moves the inside number and none of them
touches the outside number, which is zero throughout. Netgen matches
the layout to the netlist uniquely with those macros abstracted.
What we do have is a design whose redundancy claim was checked at three
levels and was wrong at two of them.
At the register-transfer level the claim was true. In the mapped netlist
it was false---the synthesiser's structural hashing had collapsed three
declared replicas into one physical bank at the geometry then being
hardened, and the register-transfer-level
fault-injection campaign that was supposed to detect this produced
\emph{bit-identical} output on the defective and the repaired design,
because it deposits into signals the netlist no longer contains.
Repaired, and confirmed by a per-instance-path census of mapped
flip-flops that is mutation-checked and runs in two independent
technology mappers.
In the die the three banks are triplicated but not physically
independent, and we publish the measurement with the instrument that
makes it re-runnable: the banks share one region of the layout, the
minimum cross-replica separation is \SI{3.78}{\micro\meter}
centre-to-centre, which is the floor for two cells in adjacent rows, and
49 pairs abut. Against a null model that shuffles only the replica
labels over the same placed positions, the banks are statistically
indistinguishable from an arbitrary labelling; what \emph{is} measured is
a \textbf{4.1-fold concentration on corresponding bits}---median
\SI{27.52}{\micro\meter} against \SI{113.47}{\micro\meter} over all
cross-replica pairs---which is the shape a shared voter's wirelength
would leave.
Nothing in the flow asked for separation, and the census that establishes
the redundancy runs on the synthesis netlist and cannot see placement by
construction.
A third result is orthogonal and sharper: in a gate-level campaign on the
watchdog's protection banks, 174 of 174 injected upsets classified
\textsc{corrected}---and every one of them reset the part one clock
later, through a combinational decode of the voted state onto an
asynchronous reset that the register-transfer level cannot express.
The generalisable claim is not a taxonomy of defects. It is that a
fault-tolerance property is not one claim but a claim at each level
between the source and the die, that the instrument which establishes it
at one level is structurally incapable of refuting it at the next, and
that each replacement instrument must carry a mutation that has been seen
to make it fail. We publish the instruments, the mutations, the negative
results and the claims we withdrew.
\end{abstract}

\section{Introduction}
\label{sec:intro}

This paper is about a claim we made about our own design and could not
support, twice, at two different levels, while the checks we had pointed
at it all passed.

The design is a radiation-tolerance-oriented neuromorphic
system-on-chip built entirely with open tools on an open 130\,nm PDK,
through a community-governed successor of the tapeout-proven open
register-transfer-to-mask-data flow~\cite{openlane2020}.
Its configuration state is triplicated and majority-voted; its memories
and register file carry SEC-DED codes; its watchdog is windowed and its
protection state is triplicated with a storage transform. All of that is
ordinary, and none of it is this paper's contribution.
The contribution is what happened when we tried to establish that the
redundancy was actually there.

\paragraph{The first level.}
The register-transfer source declares three replicas. That is checkable
by reading it, and it was true.

\paragraph{The second level.}
The mapped netlist did not contain three replicas. Yosys's
\texttt{opt\_dff} normalises hold-multiplexers into enable flip-flops and
\texttt{opt\_merge} then hashes cells on their (D, EN, reset-value)
signature; three replicas written from the same expression hash
identically and become one. The netlist that fed the \textbf{4$\times$2}
hardening run contains 362 references to \texttt{cfg\_a}, none to
\texttt{cfg\_b} and none to \texttt{cfg\_c}; at that geometry the RTL
declares \num{1161} flip-flops and the netlist has \num{1045}. The
geometry is printed because Section~\ref{sec:design} states the rule
that it must be---the frozen design, the placement measurement and the
55/55/55 census are all 6$\times$2---and an earlier draft of this paper
broke that rule five times in the document that states it.

This is not a new phenomenon. It is documented in the FPGA
radiation-hardening literature~\cite{kastensmidt2005}, in vendor
guidance~\cite{esa_ftmr}, and it is the reason post-synthesis TMR
insertion tools exist~\cite{tmrg2017,johnson2010}. What is worth
reporting is not that it happened but \emph{what failed to notice}: a
committed register-transfer-level fault-injection campaign, aimed
specifically at the configuration TMR, reporting five injections per
replica, all \textsc{corrected}, zero silent data corruptions---and
producing the same numbers before and after the repair, because at the
register-transfer level the three replicas are three distinct signals
regardless of what the netlist holds. A second check, an A/B of the total
flip-flop count with and without a preservation attribute, returned the
same total both ways and that identity was read as ``nothing was
merged''. It is equally consistent with ``the merge happened and the
attribute does not prevent it'', which is what was true: the attribute
lands on the net, not the cell, so 110 names survived as aliases while
the storage was deleted.

\paragraph{The third level.}
The repair was made, and the instrument built for it---a census of mapped
flip-flops by \emph{instance path}, in two independent technology
mappers, with the preservation attributes deleted from the source text,
carrying a mutation that makes it fail---has passed ever since. We then
asked a question nothing in the project had asked: \emph{where are the
three banks on the die?}

They are in the same place. Section~\ref{sec:placement} is the
measurement. The census establishes redundancy in the netlist and is
structurally incapable of seeing placement, because it runs before
placement exists. A guard must sit at the stage that can satisfy it, and
this one cannot.

\paragraph{What this paper claims, and what it does not.}
It does not claim the design is radiation-hardened. It cannot: there is
no silicon, no beam, no total-ionising-dose curve, no single-event
cross-section, and the design maps to stock standard cells rather than to
a radiation-hardened-by-design library, one of which exists for a
neighbouring PDK and is not used here~\cite{sky130rhbd}. It does not
claim a failure rate; every fault-injection figure is conditional on an
upset having landed in the injection window. It does not offer a
cost-versus-residual-SDC curve and therefore makes no cost-optimality
claim, which is the axis CLEAR~\cite{clear2018} settled.
It claims exactly this: the redundancy is real in the netlist, is not
real in the die, and each of those statements needed a different
instrument to establish.

\paragraph{Structure.}
Section~\ref{sec:inventory} is an inventory of what exists and what does
not, before any argument, because the single largest risk to a paper like
this is a reader discovering the absence of silicon halfway through.
Section~\ref{sec:design} describes the design;
Section~\ref{sec:mechanisms} the protection mechanisms and their measured
cost; Section~\ref{sec:formal} what is proved and what is assumed;
Sections~\ref{sec:fi-method} and~\ref{sec:fi-results} the fault-injection
method and results; Section~\ref{sec:placement} the placement
measurement; Section~\ref{sec:evidence} the paper's central argument;
Section~\ref{sec:determinism} the flow determinism that licenses every
delta in the paper to be attributed at all; Section~\ref{sec:limits} the
limitations; Section~\ref{sec:related} related work, which credits away
most of what this design does.

\section{What exists and what does not}
\label{sec:inventory}

\paragraph{Every layout-derived number in this paper predates a
register-transfer change made after it was written, and that is stated
here rather than discovered.}
On 11 September 2026, five defects were repaired across three
system-on-chip sources---the accelerator, the watchdog and the console
transmitter: a deadlock in the event engine's decision state from which
it could not drain the die's queue, an asynchronous board strap sampled
into three triplicated replicas with no synchroniser, a slave that
advertised itself free while it still owed the fabric a response, three
counter comparisons written as equality where the file's own rule is
inequality, and a holding-register race that silently swallowed one
transmitted byte in eight. All five are real, all five are repaired, and
each carries a test that fails on the unrepaired source.

\emph{Corrected 12 September 2026: this paragraph said four defects in
two files. The fifth is the console transmitter's, and it was the one
most easily missed because it is reachable only by a driver that writes
without polling---which no program in the repository does, so no
simulation had ever produced it. Omitting it from a marker about
completeness is the failure the marker exists to prevent.}

\textbf{None of the layouts this paper measures contains them.} Every
\file{soc\_top} run cited here---the inventory of
Table~\ref{tab:inventory}, the deck results, the $2\times2$ timing
factorial, the corrected-upset campaign of Section~\ref{sec:w9} and the
gate-level divergence of Section~\ref{sec:fi-results}---was hardened
before the repair, and the repair adds five flip-flops and about
twenty-three cells. \emph{Those numbers remain true of those runs}, which
is what the digest ledger pins and what
Section~\ref{sec:determinism}'s determinism licenses. What is no longer
true is the implicit reading that they describe the design as it now
stands in the repository.

We report it this way, rather than re-hardening, for a stated reason: a
re-harden plus a re-run of the deck, LVS and static-timing evidence is
seventeen hours for the Magic deck alone, and it would replace a set of
numbers whose provenance is complete with a set measured in a hurry
against a deadline. The honest alternative is the marker. What it costs
is that a reader who clones the repository and hardens it will not
reproduce the figures below, and the difference is five flip-flops.

\textbf{And the reason this is printed at the top of the inventory rather
than in the limitations is the shape of how it was nearly missed.} On the
same day the design moved, one instrument pinned 156 new digest rows
attesting exactly the superseded evidence and another made this paper's
numbers about that evidence more precise. Both were correct---identity is
identity, and 1{,}806 really is the right pin count---and the effect of
the two together is that every check in the repository now reports green
over evidence that has quietly stopped describing its subject. That is
this paper's own subject, arriving through the door marked \emph{success}
rather than \emph{failure}, and no single check could have seen it.

Table~\ref{tab:inventory} is the inventory, and the rows saying
\emph{does not exist} are the reason it is on the second page rather than
in a limitations section at the end.

\begin{table}[htbp]
\centering
\caption{What has been built and measured, and what has not.
Figures registered in \file{paper/claims.yaml} are checked by
\file{paper/check\_claims.py}; the rest are read from the run reports and
documents cited in the text and are not machine-checked.}
\label{tab:inventory}
\small
\begin{tabularx}{\textwidth}{L{4.6cm}X}
\toprule
\multicolumn{2}{l}{\textbf{The pilot: a 12-tile shuttle submission, frozen by git blob hash}}\\
\midrule
Geometry & 6\,$\times$\,2 tiles; die \SI{404499}{\micro\meter\squared},
core \SI{392988}{\micro\meter\squared} \\
Content & 8 neurons, 8 axons, flip-flop storage, 4-deep event queues \\
Mapped cells & \num{12546} excluding fill; \num{1296} flip-flops \\
Geometric sign-off & Magic DRC 0, KLayout DRC 0, KLayout XOR 0, Netgen LVS
0 on all four unmatched counters, antenna 0, route DRC 0 \\
Timing & slow-corner setup \textbf{+\SI{1.2262}{\nano\second}} at a
\SI{20}{\nano\second} period with a real 5\,\% on-chip-variation derate;
0 setup and 0 hold violations at all three corners \\
\midrule
\multicolumn{2}{l}{\textbf{The system-on-chip around it}}\\
\midrule
Geometry & \SI{2306.40}{\micro\meter} $\times$ \SI{2075.22}{\micro\meter}
= \SI{4.7863}{\milli\meter\squared} \\
Mapped cells & \num{176447} instances including fill and six SRAM
macros; \num{5872} flip-flops --- run \texttt{s77gate}, the same run the
timing below is measured on \\
Timing & \textbf{$-$\SI{4.6799}{\nano\second} slow-corner setup on
\num{3088} violating endpoints}; hold met at all three corners;
detailed-route DRC 0 \\
Geometric decks & Run for this paper, having been off in the
configuration since the design existed. Magic DRC \num{39969214} error
boxes, which the report itself says to divide by 3 or 4, and KLayout
\num{9668} markers, of which \num{3918} are a second rule repeating the
first rule's shapes, so \num{5750} distinct; \textbf{all of both inside
the six vendor SRAM macros and none outside}, on every one of those
divisors; Netgen LVS \emph{circuits match uniquely}
with those macros abstracted to their pin lists. \S\ref{sec:decks} \\
Not covered by that & metal density and antenna are separate runsets and
were not invoked; LVS excludes the macro internals by construction. The
vendor macro itself fails at a rate linear in its bit count and cannot be
fixed here \\
\midrule
\multicolumn{2}{l}{\textbf{What does not exist}}\\
\midrule
Silicon & none, no engineering sample, no schedule for one \\
Radiation data & none we could find. We located no published TID, SEE or
SER characterisation for either PDK, and neither ships such data in its
own release \\
Application accuracy & no classification result for this design on any
task, and none has ever existed \\
A clock frequency for the SoC & the \SI{20}{\nano\second} target is not
met and a measured non-closure is not a specification. What \emph{is}
measurable is the period at which this layout closes,
\SI{25}{\nano\second}, and \S\ref{sec:decks} prints it as the property of
one layout rather than as a specification of a design \\
\bottomrule
\end{tabularx}
\end{table}

Two of those rows deserve a sentence each rather than a footnote.

\paragraph{The SRAM macro, and a decision that was narrower than it
looked.}
Six macros were needed and one was integrated as a bring-up vehicle. Its
vendor views fail three independent decks by margins that are not
marginal: \num{1106478} Magic errors, \num{2316} KLayout errors
attributed to 57 cells all inside the macro hierarchy, and 359 Netgen LVS
errors. On that evidence the three decks were switched off in the
system-on-chip's configuration, and they have never been run on the
system-on-chip at all.

Two things about that are worth separating, because the project's own
record ran them together.

The \emph{diagnosis} is well supported and we stand behind it. All
\num{1106478} Magic errors fall inside the macro footprint and
\textbf{none falls outside}; the KLayout errors all name cells inside the
macro hierarchy; and the LVS failure is not geometric at all but a
\emph{naming} mismatch---the mask-data cell names and the circuit
netlist's subcircuit names disagree and the bus delimiters differ, so
none of the 188 bus pins match. That is an upstream integration failure
with an upstream owner.

The \emph{decision} was described in the project's record as a choice
between running the decks and substituting the PDK's own unverified
residual rule set, with the substitution examined and declined on the
grounds that a sign-off obtained by relaxing a deck is not a sign-off.
That reasoning is sound and the option space was not. \textbf{A third
route existed and was not examined}: abstracting the vendor macro so that
the deck reads its boundary rather than descending into its geometry,
which is what an integrator ordinarily does with a hard macro whose
internals the vendor certifies. It is not deck relaxation---it changes what is looked at, not what
counts as an error---and the diagnosis above is precisely the evidence
that it would work, because a deck that never enters the macro cannot
report errors that are all inside it.

\textbf{That reasoning is sound and, for one of the two decks, it is
also wrong}, which we could only learn by running both.
Section~\ref{sec:decks} reports the measurement: abstracting the macro
for Magic's deck \emph{introduces} 123 violations that do not exist in
the mask data, because the tool paints the macro's boundary obstructions
as though they were drawn metal and four wide-metal rules then fire on
them. So the option that was never examined turns out to be the wrong
one here---and the seventeen-hour run over real geometry, which is what
we did instead, is the one that answers the question. The general
argument survives; its application to this deck does not.

We record this as an instance of the paper's own subject arriving one
level up from a check: \emph{the option space was described as explored
when a standard option in it had not been looked at}. The measurement
that follows from taking that route is reported in
Section~\ref{sec:decks}.

\paragraph{The frequency.}
The project operates a standing rule that a measured non-closure is not a
specification. The SoC misses its target; publishing the number it misses
by, as though it were an achieved frequency, is the class of overclaim
this paper is about. The pilot's number is published because the pilot
closes.

\section{The design under study}
\label{sec:design}

\subsection{Two designs, and the gap between them}

The specification defines a 512-neuron, 512-axon leaky-integrate-and-fire
core with 4-bit signed weights in an axon-major crossbar, \num{256}k
synapses in SEC-DED-protected SRAM, and ten bit-exact integer equations
that a Python golden model implements as the executable specification.

\textbf{That core has never been elaborated.} The top-level module
refuses geometries outside $[4,16]$, and the fallback trigger that
governs this---storage moves from SRAM macros to flip-flops when the
macros are unavailable---has fired, because the macros are the ones
Section~\ref{sec:inventory} describes. What was built, frozen and
submitted is 8 neurons and 8 axons with flip-flop storage: the
specified \emph{datapath} at a fallback \emph{scale}.

Every result in this paper is a result about the built design. Where a
figure is measured at a particular geometry, the geometry is printed with
it. Nothing here is a claim about the specified core.

\subsection{The reference, and what is and is not claimed about it}

The accelerator's architecture takes its shape from a published product
brief for a 28\,nm commercial part, of which two pages are public and no
datasheet, user manual or conference paper exists. We use it as an
\emph{architecture-class reference} and never as a performance baseline.
No performance comparison is offered, and the one figure this project
once derived from that comparison has been withdrawn: the source list
mixed items from the two-page brief with an item from the vendor's
product page under a single heading naming the brief, and the resulting
compute ratio inherited a provenance it did not have.

There is a second open question we do not resolve. Whether this
accelerator is an independent implementation or a reference-guided one
has not been settled in the project's own record, and until it is,
\textbf{no claim of independent implementation is made here}. The word
that would usually be used does not appear in this paper.

\section{Protection mechanisms and their measured cost}
\label{sec:mechanisms}

\subsection{A replication bound that comes from the tool, not from
information theory}

Three replicas of a $W$-bit word, written from the same expression, are
merged by structural hashing. The project's response is to give each
replica a different \emph{storage transform}: replica A stores the value,
B stores it under a polarity mask, C under a mixing transform composed
with a mask.

The available affine transforms over $W$ bits give exactly
\[
2\,(2^{W}-1)
\]
distinct coordinate functions---2 at $W=1$, 6 at $W=2$, 14 at $W=3$.

Three replicas of a $W$-bit word need $3W$ pairwise-distinct coordinate
functions, which at $W=2$ is exactly the six available, \textbf{so
counting alone does not settle $W=2$}: the assignment
$A=(x_1,x_2)$, $B=(\bar x_1, x_1\!\oplus\!x_2)$,
$C=(\bar x_2, \overline{x_1\!\oplus\!x_2})$ gives three distinct
invertible replicas at two bits. What settles it is the construction
this design uses. With $A$ on the identity and \emph{$B$ pinned to the
polarity transform}, the two of them consume
$\{x_1,x_2,\bar x_1,\bar x_2\}$, and the two functions that remain are
complements of one another---so a third replica built from them would
store one bit twice and determine nothing. \textbf{Under this transform
family, with $B$ on polarity, three bits is the first width at which a
third replica has functions left to take.}

The implementation is narrower still, and says so at elaboration:
\file{hw/soc/rtl/soc\_tmr\_bank.v} refuses $W<4$, because below four
bits the mixing layer's rotation degenerates and a stored row can fall
back to weight one. A reader following the counting rule alone would
reach an elaboration error in this repository, which is why the rule is
printed with its restriction rather than as a bound.

This is a statement about defeating a specific optimisation over a
specific family of transforms, under a specific choice for the second
replica---not an information-theoretic bound, and not a general claim
about three affine replicas of a two-bit word, which exist. It should be
read that way, and an earlier draft of this paper did not: it printed the
$W=3$ conclusion directly under an arithmetic that makes $W=2$
sufficient, with ``one of the three is the identity'' standing in for the
argument. That is the paper's own subject in the paper's own contribution,
and it was caught in review rather than in writing. It is nonetheless a design rule with immediate
consequences that are visible in the source: a one-bit flag gets two
rails and detection only, never a triple; a two-bit state gets a check
field rather than a third copy; and scattered one-bit flags are
\emph{bundled} into words of 22 or 29 bits (one bank in its two
configurations), 23 and 25 before triplication, because a bundle has
functions to spare and a flag does not. Three structures; an earlier
draft listed four by counting one bank's two configurations separately.
We found no statement of this bound in the surveyed literature, though
the underlying problem is thoroughly documented~\cite{kastensmidt2005,
esa_ftmr,tmrg2017}.

\subsection{Codes}

The weight path uses a (72,64) Hsiao SEC-DED code~\cite{hsiao1970}. Its
parity matrix was re-derived from the source constants rather than taken
on trust: all 72 columns are distinct, non-zero and of odd weight---56 of
weight 3, 8 of weight 5, 8 identity---and every row has weight 26 over
the data field, so all eight parity trees are five levels of XOR2. Odd
column weight is the double-error-detection argument: a two-bit error
always produces a syndrome of even parity, which can never equal a
column.

Shortened to 32 data bits one row vanishes, giving (39,32) and seven
check flip-flops per register---the minimal SEC-DED width for 32 bits.
Shortened to 8 it becomes (13,8) inside a 16-bit lane, and four
independent lanes fill a 64-bit memory row exactly. That last choice has
a consequence worth stating because it is the reason the memory was
halved rather than widened: two upsets in two lanes are two corrected
singles, where a single word-wide code would report one uncorrectable
double.

Formally proving an ECC decoder is routine industrial
practice~\cite{kumar2023} and is not claimed here as a contribution.

\subsection{Cost, and the rule that makes a cost figure mean anything}

Costs are measured against \emph{the same design with hardening disabled},
never against the earlier design. The difference is not pedantic. When a
hardening change is accompanied by a refactor---and it usually is---
measuring against the earlier design credits the refactor to the
hardening. The project records five instances where this was caught after
the fact, and \textbf{four of the five were in the direction that
understates the cost}: \SI{683.1216}{\micro\meter\squared},
\SI{58.17}{\micro\meter\squared}, \SI{27.216}{\micro\meter\squared} and
\SI{25.4394}{\micro\meter\squared}. The fifth,
\SI{10.9242}{\micro\meter\squared}, runs the other way, and the document
that found it says so in bold. An earlier draft of this paragraph printed
the universal --- in the subsection about the rule that makes a cost
figure mean anything.

The register-file SEC-DED substitution costs
\SI{38050.5762}{\micro\meter\squared} (+13.80\,\%) and 222 flip-flops on
the core, measured module-alone in one session. Two qualifications travel
with that figure and are printed here because they are printed in the
record: an earlier section of the same document states
\SI{37983}{\micro\meter\squared} for the same measurement, a
\SI{67.58}{\micro\meter\squared} discrepancy that stood uncorrected until
this paper was prepared; and the module-alone configuration is not the
shipped one, which measures \SI{316051.6590}{\micro\meter\squared}
(+14.64\,\%).

\section{Formal verification: what is proved, what is assumed}
\label{sec:formal}

\subsection{The scope, counted from the files}

There are 23 SymbiYosys job files defining \textbf{118 tasks}: 54 over
the frozen pilot and 64 over the system-on-chip fabric and peripherals.
By mode, \textbf{51 are \texttt{prove}} by $k$-induction, 36 are bounded
model checking and 31 are cover, with depths from 1 to 220.
\textbf{Mode is not the same axis as boundedness}, and an earlier draft
of this paper conflated them by calling the $k$-induction results ``the
only unbounded results anywhere in the project''. Ten \texttt{bmc} tasks
are \emph{exhaustive} rather than bounded because their models are
stateless: the two SEC-DED codecs, the LIF codec and the voter at three
widths run at depth~1 over free symbolic inputs, so one step ranges over
the whole input space---and the voter's own job file says so in its
header. The equivalence proofs of Section~\ref{sec:evidence} are
unbounded for the same reason. All 118 have a
result directory and all 118 report \texttt{DONE (PASS)}. The task set
and the result-directory set are exactly equal: 118 expected, 118
present, none missing, none orphaned.

\subsection{Three things about that result that a reader must have}

\paragraph{No verdict is committed.}
\texttt{git ls-files '*logfile.txt'} returns \textbf{zero}. Every result
directory is build output and is ignored. The only tracked artefact is a
verification log that records \emph{counts}. A reader who clones this
repository obtains 23 job files and no verdicts, and must re-run to see a
single \texttt{PASS}. We state this rather than write ``118/118 pass''
unqualified.

\paragraph{The date of the evidence is not in the evidence.}
SymbiYosys log lines carry a wall-clock time of day and no date. The
figure this project reports as the age of its oldest verdict is derived
from file modification time, which a clone, a checkout or a \texttt{touch}
resets and nothing recovers.

\paragraph{Whether the verdicts match the current source is checkable,
and nothing in the project was checking it.}
A SymbiYosys run copies its sources into the task directory. Comparing
all 351 such copies against their tracked originals---336 \texttt{.v} and
15 \texttt{.vh}, the header files being the class where drift is
likeliest and the class an earlier draft silently dropped---raw
comparison says they differ, and the entire difference is a three-line
licence header added later. Stripping only those headers gives
\textbf{351 identical, no real difference}, so the verdicts do correspond to the design as it
stands. This was established once, by hand, for this paper, and it is now a
check.

The project's own verification recorder had two blind spots and they are
worth reporting because the file exists to stop exactly this. It read the
\texttt{DONE} line and the file's timestamp and nothing else, so a log
left over from superseded RTL would have been counted as a pass. And it
enumerated results with a two-level wildcard over two directories, which
\textbf{structurally could not reach} the processor checks of
\S\ref{sec:rvformal} at all: it reported \num{118} passing and
\emph{zero} not-passing on a tree that also held 6 failures, 3 errors and
26 stopped runs. Every row it had ever written said
\texttt{formal\_other=0}, and the statement was true of what it looked at
and false of the repository.

Widened to every directory holding both a job file and a log, and taught
to compare each task's own source copies against the tracked originals
modulo comment-only differences, it now reads \num{425} passing and
\num{35} not, over 460 result directories, with no stale verdict and no
unchecked one. \textbf{The numbers in this section did not change. The
recorder's ability to contradict them did.}

\subsection{The assumption ledger}

Across 25 property files there are 106 \texttt{assume} statements, 15
\texttt{\$anyconst} symbolic constants and 9 \texttt{\$anyseq} symbolic
sequences, and \textbf{zero} \texttt{restrict}. Ten further
\texttt{assume} statements live in the processor harness wrapper, and
they are the ones Section~\ref{sec:rvformal} criticises---an earlier
draft's ledger excluded them and then attacked two of them.

\emph{Five} property files are the top of their own proof and carry
\textbf{no} assumptions at all: the SEC-DED codec, the byte-lane codec,
the register-file codec, the LIF memory codec and the TMR voter. These
are the strongest results in the project, and the SEC-DED file states
each property as an implication on Hamming weight rather than assuming
one. A sixth, the memory-ECC integration set, carries no assumption of
its own and is \texttt{`include}d into a model that assumes reset at
time zero and no request while in reset; it is counted here as assumed,
because \textbf{an assumption count per \emph{file} is not an assumption
count per \emph{proof}}---and counting it at the wrong level returned a
green result inside the ledger built to prevent that.

The load-bearing assumption is narrower than an earlier draft claimed.
\textbf{Every fault-tolerance proof about a \emph{voted} structure is
scoped to at most one bad replica}: the injection vectors in the FIFO,
configuration-voter and watchdog property sets are constrained to one
replica at a time. Nothing below that is assumed---the voter's masking
theorem quantifies over \emph{every} value of the third replica, so it
covers a fault of any weight in one replica, and the SEC-DED proofs
drive a free 72-bit error vector and state weight as an implication,
including the two-bit case.

What no proof here can express is an upset that corrupts \emph{two of
three replicas}, and Section~\ref{sec:placement} is precisely about a
geometry in which that is the relevant question. The earlier draft said
``the single-fault assumption appears in 14 of the 25 property files'',
a count that appears in no document, script or registry in this
repository and that swept in one-hot \emph{functional} assertions which
are not fault-model assumptions---and it stood three lines below a file
with no assumptions at all whose header says it covers every fault of
every weight.

\subsection{The processor: 79 checks, none of them unbounded}
\label{sec:rvformal}

The core is a small RV32IMC implementation with physical memory
protection, with its architectural register file substituted for a
SEC-DED-protected one. It was run against the standard RISC-V formal
verification harness in two configurations, upstream and substituted.

\begin{table}[htbp]
\centering
\caption{RISC-V formal checks, both configurations, parsed from every
generated \texttt{bmc} task's own log. The same configuration also
generates 79 \texttt{cover} jobs per run into a separate directory; all
79 pass in both and they are not tabulated here.}
\label{tab:rvformal}
\small
\begin{tabular}{lrr}
\toprule
& upstream register file & SEC-DED register file \\
\midrule
\texttt{PASS} & 70 & 69 \\
\texttt{FAIL} & 3 & 3 \\
no \texttt{DONE} line (stopped) & 6 & 6 \\
\texttt{DONE (ERROR)} & --- & 1 (\texttt{reg\_ch0}) \\
\midrule
generated check directories & 79 & 79 \\
of which \texttt{mode prove} & 0 & 0 \\
\bottomrule
\end{tabular}
\end{table}

Three points, in the order that matters.

\paragraph{Zero of the 79 are unbounded.}
78 are bounded model checking and one is cover. Nothing about this core
is proved for all time by this harness; it is checked to a depth.

\paragraph{The failing and stopped sets are identical across the two
configurations}---\texttt{insn\_div}, \texttt{insn\_rem} and
\texttt{liveness} fail in both, and the six multiply, unsigned-divide and
unsigned-remainder checks stop in both. That identity is the substance of
the claim that the substitution changed nothing, and it holds.
\texttt{liveness} is a statement about the bound rather than about the
core: a 37-cycle divide has not retired by the check's cycle 50.

\paragraph{The one check that behaved differently did not time out, and
the project said it did.}
\texttt{reg\_ch0} is the check over the substituted register file. Its
own log ends \texttt{engine\_0: <SIGTERM>}, ``Engine terminated without
status'', \texttt{DONE (ERROR, rc=16)} at 1\,h\,41\,m. That is an
external kill, not a bound reached. The project's summary table printed
\textsc{timeout}; its own cost ledger, four sections later, said ``killed
from outside''. The two disagree and the log settles it. Elsewhere---an
engine sweep run under \texttt{timeout(1)}, where the exit status is
\texttt{rc=124}---\textsc{timeout} \emph{is} the right word and is
verified.

We report this at length for a reason that is not pedantry.
\textsc{timeout}, \textsc{killed} and \textsc{error} are three different
words describing three different states of knowledge: a bound was
reached, something outside intervened, or the tool failed. Conflating
them converts ``we do not know'' into ``we ran out of time'', which reads
as a resource problem rather than an open question. \textbf{Nothing here
is evidence that the substitution preserves the property, and nothing
here is evidence that it breaks it.}

\paragraph{What the harness does not reach at all.}
Zero control-and-status-register checks are generated, because the core's
formal interface exposes no CSR ports. Memory-protection enforcement is
unproved, and two assumptions exist to keep it from biting. Eight of the
70 instruction models have no verdict in either run, and they are
\emph{all eight} multiply-divide instructions---25 of the 70 are
compressed, and none of the eight is base-integer. Bus errors are
assumed away while the fabric raises one on every unmapped address. Trap
destination and the three trap CSRs are unchecked.

\subsection{A wrong specification manufactures a bug}

Two of the three failing checks are not failures of this design. The
harness's own reference models for signed division and remainder compute
the \emph{unsigned} result, confirmed on four counterexamples from two
independent runs and checkable by hand from the language standard's
sign-extension rules. Corrected copies are carried in this repository
with the upstream licence notice inline.

Running them is the second half of the result, and it is a negative one:
\textbf{neither closed}. The two checks were stopped without a verdict
after 40\,m\,43\,s and 48\,m\,52\,s. Against the \emph{defective} model
the checks failed quickly, because exhibiting one counterexample is a
satisfiability question; against the \emph{corrected} model nothing was
proved. \textbf{Correcting a specification made the proof harder rather
than easier}, and the consequence is that all eight instructions of the
multiply-divide extension have no formal result in this project at all.
A paper whose stated policy is to publish negative results stopped one
sentence before its own.

The general point is the sharpest available demonstration of a thing this
paper keeps returning to: \textbf{a trusted specification is a trusted
component}, and a wrong one does not merely fail to catch a bug, it
manufactures one. Had these two checks been believed, the design would
have been ``repaired'' to compute an unsigned division.

\section{Fault injection: method}
\label{sec:fi-method}

\subsection{Fault model, stated before any result}

Single-bit, single-event, flip-flop only. One bit of one sequential
element is inverted once, at one clock edge, and the simulation continues
to completion against a golden run. There are no single-event transients,
no multiple-bit upsets, no permanent faults and no spatial correlation
between injection sites.

Two consequences follow and are stated here rather than in the
limitations, because they are large.

First, \textbf{a combinational decoder is not a site}. The register
file's SEC-DED decoder is 528 combinational cells and none of them can be
injected by this method, so the campaign says nothing about a transient
in the decoder that produces a wrong correction.

Second, \textbf{the hazard that Section~\ref{sec:fi-results} repairs is
precisely a transient}, and the campaign that measures the repair injects
only into flip-flops. That is not a contradiction---the repair is
observed through a flip-flop upset that reaches a combinational
decode---but it is a limit, and the geometry of
Section~\ref{sec:placement} is exactly where multiple-bit upsets would
matter and this model cannot express them.

\subsection{Outcome classes, and how they map onto the standard ones}

Each injection is classified into one of five: \textsc{masked} (the
upset never reaches an observable), \textsc{corrected} (a protection
mechanism fired and the answer is right), \textsc{detected} (an error was
reported and the answer is not silently wrong), \textsc{sdc} (the answer
is wrong and nothing said so), \textsc{hang} (the run did not terminate).

\textsc{sdc} and \textsc{detected} are Mukherjee's silent data corruption
and detected unrecoverable error~\cite{mukherjee2003,mukherjee2008}, and
\textsc{masked} is their benign complement, which
Shivakumar~\emph{et~al.}~\cite{shivakumar2002} decompose into logical,
electrical and latching-window masking---of which this method observes
only the first, because it injects into state rather than into a node.
The campaign structure follows Cho~\emph{et~al.}~\cite{cho2013}, which is
also the citation for what a higher abstraction level costs, and is the
reason the campaigns here are run at gate level as well as at
register-transfer level.

\textsc{corrected} is the class that is not in those taxonomies. It is
the case where a mechanism fired and the answer is right, which the
architectural-vulnerability vocabulary folds into benign. Keeping it
separate is what makes Section~\ref{sec:fi-results}'s central result
expressible at all: a class where everything went right by every measure
the classifier had, and the machine restarted.

\subsection{A weakness we did not repair}

We report proportions and, in most places, do not report confidence
intervals derived from statistical fault-injection
theory~\cite{leveugle2009}. Where the population is small the project
uses the rule of three---``0 of 100 is a rate below about 3.7\,\%, not a
rate of zero''---and that phrasing appears in the record. It is not
enough, and for a paper about evidence discipline it is the most
uncomfortable gap in this section. It is a day of work and it was not
done before this submission.

\subsection{The bench was wrong, and the check that caught it}

The system-on-chip's injection bench decides its deposit instant by a
wait loop racing its own cycle counter. Over a 1{,}400-injection
campaign, \textbf{686 deposits landed one clock edge after the cycle
their record names}, and 714 landed on it.

This bench passed its own per-deposit verification on all
1{,}400 injections of that campaign (2{,}800 simulations, golden and
injected)---each injection is checked to have flipped the bit it
names---and it passed a uniquely-determined cross-design trace alignment.
Both are real checks and neither can see the defect, because both ask
\emph{was the right bit flipped} and the defect is \emph{when}.

What caught it was a different question: a cycle-for-cycle comparison of
every mapped flip-flop between the register-transfer and gate-level
models of the same record. On the first record whose outcome depended on
the edge, the two levels disagreed completely---an upset that derailed
the register-transfer model with 64 traps and a nonzero exit status was,
one edge earlier, \emph{a golden run}. Before the repair that control
passed on three records and failed outright on the fourth; after, it
reports \textbf{0 of 964 differing flip-flops over 18{,}681 cycles on all
four}.

The consequence was not a patch. A 450-record gate-level pass, seven
hours and twenty-two minutes of compute, was \textbf{discarded} rather
than corrected, because a campaign whose deposit instant is not known is
not a campaign with a small error. The published outcome distributions
did not change; what changed is that they are now attributable to an
instant that was measured.

\section{Fault injection: results}
\label{sec:fi-results}

\subsection{Register-transfer and gate level do not diverge much, and
where they do the mechanism is nameable}

On the frozen pilot, 370 of 370 like-for-like injections classify
identically between the register-transfer model and the sign-off netlist.
On the system-on-chip, of 1{,}175 records replayed at gate level,
\textbf{1{,}158 classify identically and 17 disagree}. Both disagreement
structures were characterised rather than reported as noise: fifteen are
injections into an instruction register that the synthesiser stored once
for two register-transfer copies, so the netlist is strictly worse on all
fifteen; two are the multiplier's re-encoded state, where the netlist is
\emph{better}---silent corruption at register-transfer level, detected on
the netlist. Separately, 75 injections went into netlist flip-flops with
no register-transfer twin, the most lethal storage in the design per bit,
and those have no register-transfer record to disagree with at all.

We draw the conclusion the numbers support and no more: at this scale,
on this design, register-transfer-level injection is a good
approximation of gate-level injection for \emph{outcome distribution},
and it is not a substitute for it, because the cases where it differs are
structural and are invisible at the level that approximates well.

\subsection{A correction that cost a reset}
\label{sec:w9}

This is the sharpest single measurement in the project.

The watchdog's protected state is 29 bits, triplicated into three banks
of 29 flip-flops. A gate-level campaign injected \textbf{174 single-bit
upsets}---87 replica flip-flops, two draws each, 58 per bank---into the
sign-off netlist. Every one classified \textsc{corrected}. No error was
announced. Masking fields were identical on every record: nothing
\textsc{masked}, nothing \textsc{sdc}, nothing \textsc{detected}, nothing
hung, the voter's mismatch flag set once, the output correct.

An \emph{event} counter on the reset pin---a channel distinct from the
clocked sampler the classifier reads---recorded that \textbf{every one of
those 174 corrections also asserted the system reset}, one clock later.

\begin{table}[htbp]
\centering
\caption{One byte-identical 174-injection plan, four netlists.
Every masking field is identical across all four arms; only the reset
event count moves.}
\label{tab:w9}
\small
\begin{tabular}{llrrr}
\toprule
Design & Netlist & \textsc{corrected} & reset events & resets/correction \\
\midrule
unrepaired & placed      & 174/174 & 174 & 1.0000 \\
repaired   & placed      & 174/174 & 0   & 0.0000 \\
unrepaired & synthesised & 174/174 & 174 & 1.0000 \\
repaired   & synthesised & 174/174 & 0   & 0.0000 \\
\bottomrule
\end{tabular}
\end{table}

\paragraph{The mechanism.}
The voted 5-bit reset-hold field was decoded \emph{combinationally} onto
the asynchronous reset of every flip-flop in the system-on-chip. A
correction is a state change; the decode of that change reached the
reset before it was registered. In the source this is a single
continuous assignment followed by an \texttt{always} block sensitive to
the resulting net, under a comment asserting that the request is
registered.

\paragraph{Why nothing saw it.}
The gate-level classifier samples announcement pins on the rising clock
edge and cannot see a sub-clock pulse on an asynchronous reset. Every
register-transfer-level campaign missed it---126 and 168 injections, all
corrected, no reset---and so did the formal property, because
\texttt{rst\_req\_o = in\_reset} and
\texttt{rst\_req\_o = in\_reset \&\& in\_reset\_q} are \emph{the same
function of the same state at the register-transfer level}. The property
passed by $k$-induction on the defective design and on the repaired one,
and it was true of both.

There was one further tell that was visible and not looked at: on all 174
records of both unrepaired arms, the run length minus the injection edge
is \emph{exactly} 18{,}683, with no spread. A machine that restarts
finishes at a fixed offset from the restart. A spread of zero over 174
records is not a plausible property of an unperturbed run, and nobody
asked.

\paragraph{The repair, and the instrument.}
One flip-flop: 131 to 132 in the block, 5{,}873 to 5{,}874 in the
system-on-chip, \SI{643.0536}{\micro\meter\squared} of synthesised area,
\num{0.0993}\,\%. Die area unchanged.

The instrument built for it is a whole-netlist census of asynchronous
reset fan-in. On the unrepaired netlist there are five distinct
asynchronous-reset source sets; four are a port or registered signals,
and the fifth---one net, two flip-flops---has a combinational fan-in of
\textbf{25 flip-flops, every one of them a watchdog replica bit}. On the
repaired netlist the same cone has 26 sources: the same 25 plus the
register that now stands between them and the reset. The census is a
general shape rather than a special case: \emph{flag any asynchronous
reset whose fan-in is a decode rather than a register}.

\paragraph{The mutations.}
Two, both committed and both executing in the regression: the pre-repair
design, and an inverted combining term. Both fail the proof.

\section{The two checks that were switched off, and the scope they were
switched off at}
\label{sec:decks}

Sections~\ref{sec:evidence} and~\ref{sec:placement} are about instruments
that ran and could not see. This section is about two that did not run,
and it belongs in the same paper because the reason in both cases is a
\emph{scope} that was set once and then read as though it were wider.

\subsection{The geometric decks}

The frozen pilot passes every deck an open flow can run: Magic DRC 0,
KLayout DRC 0, KLayout XOR 0, Netgen LVS 0 on all four unmatched
counters, antenna 0, route DRC 0. That is the design being manufactured.

The system-on-chip has no such result, and the immediate reason is
mechanical: \texttt{RUN\_MAGIC\_DRC}, \texttt{RUN\_KLAYOUT\_DRC},
\texttt{RUN\_KLAYOUT\_XOR} and \texttt{RUN\_LVS} are all \texttt{0} in its
configuration. The flow log records four consecutive lines of the form
\emph{``Gating variable for step \ldots{} set to `False'---the step will
be skipped''}, and the manufacturability report at the end notes that
three counters were ``not reported''.

The evidence behind that setting is real and we do not dispute it. What
we dispute is the shape of the conclusion drawn from it, which
Section~\ref{sec:inventory} sets out: the diagnosis was that every Magic
error falls inside the vendor macro's footprint and none outside, and the
conclusion drawn was that the decks would tell a reader nothing---when
the diagnosis is in fact an argument that a deck which does not enter the
macro would return a result about this project's own geometry.

\paragraph{Why abstracting the macro is not relaxing the deck.}
The distinction matters and it is the whole legitimacy of the experiment.
Substituting a permissive rule set changes \emph{what counts as an
error}, and a clean result obtained that way means less than no result.
Abstracting a hard macro changes \emph{what is looked at}: the deck runs
unmodified over every polygon the integrator drew and treats the macro as
a boundary, on the standing assumption that the vendor certifies its own
cell. That assumption is exactly what the 1{,}106{,}478-to-0 inside-to-
outside split calls into question for this PDK version---and that is a
statement about the vendor's macro, which is the finding, not a
concealment of it.

\paragraph{What the decks say when they are run.}
All three were run over this design's own final mask data. The claim was
bounded in advance---at most ``no design-rule error in the geometry this
project drew'', never ``the design is DRC-clean''---and the results are
inside that bound.

\begin{table}[htbp]
\centering
\caption{The three geometric decks over run \texttt{s77gate}, unmodified,
at the pinned tool and PDK versions. The split is a geometric test of
every error coordinate against the six macro footprints taken from the
placement, not an attribution by cell name.}
\label{tab:decks}
\small
\begin{tabularx}{\textwidth}{L{2.9cm}rrX}
\toprule
Deck & inside the macros & outside & \\
\midrule
Magic DRC, full & \num{39969214} & \textbf{0} & 17\,h\,11\,m, 16\,GiB
peak. Also 0 straddling a macro edge and 0 outside the die \\
KLayout DRC & \num{9668} & \textbf{0} & 3 of 173 rules fire; every item
names a cell inside the vendor hierarchy. \num{3918} of the markers are
\texttt{Sdiod.e} repeating \texttt{Sdiod.d}'s shapes exactly, so the
distinct-marker count is \num{5750}; the outside column is 0 either way \\
Netgen LVS & --- & --- & \textbf{Circuits match uniquely}: \num{62256}
devices and \num{61912} nets on both sides, all seven counters zero \\
\bottomrule
\end{tabularx}
\end{table}

Four things make that table a measurement rather than an assertion.

\emph{The classifier was validated against a published answer before it
was pointed at anything new.} Run against the retained report of the
earlier single-macro experiment it returns \num{1106478} boxes,
\num{1106478} inside, 0 outside, and reproduces that experiment's
ten-row rule histogram line for line.

\emph{The macro was calibrated on its own.} The same deck over each
vendor macro standing alone, with no design near it, gives
\num{8808860} and \num{2210913} boxes---and both scale \emph{linearly
with bit count} off the single-macro figure ($1.9982\times$ and
$7.9611\times$), which is the signature of a per-bitcell defect rather
than anything contextual. Six instances predict \num{39657266} against
\num{39969214} measured, a surplus of \num{311948} or 0.79\,\%, and
every box of that surplus still lands inside a macro.

\emph{Neither headline is a violation count, and the two overstate for
different reasons.} Magic's figure is a box count and the report says so
itself: its last line reads \emph{``should be divided by 3 or 4''}, so
the distinct-violation figure is roughly 10.0 to 13.3 million.
KLayout's \num{9668} needs the same treatment and an earlier draft of
this paper gave it none, printing it uncaveated in the abstract, in
Table~\ref{tab:inventory} and in Table~\ref{tab:decks} while caveating
Magic's in all three: \texttt{Sdiod.d} and \texttt{Sdiod.e} are
evaluated over the same derived layer and flag the \emph{identical}
\num{3918} (cell, geometry) pairs---intersection \num{3918}, symmetric
difference zero---so one population of shapes is reported twice and the
distinct-marker count is \num{5750}. That the two rules are one
population was established on the single-macro case in the project's own
record and reproduces here on the whole system-on-chip. The same record
measured the opposite error in the same number: a deep-mode marker count
\emph{understates} the shape count, by a factor of 53 on the single
macro, and we have not measured that factor here. So the KLayout
headline is wrong in both directions at once and its true shape count is
not known to us. In every case the divisor changes the inside number and
does not touch the outside number, which is zero on all of them.

\emph{The zero is not the deck declining to look---and that is
established by a planted control for \textbf{one} of the three decks,
not for the table.} The control is KLayout's alone: the same deck, same
version, same options---deep mode, recommended rules off---on a
synthetic top cell carrying one \SI{0.10}{\micro\meter} Metal1 wire
against a \SI{0.16}{\micro\meter} minimum width and two wires
\SI{0.10}{\micro\meter} apart against a \SI{0.18}{\micro\meter} minimum
space. It returns exactly those two violations, \texttt{M1.a} and
\texttt{M1.b}, and nothing else. \textbf{No control with a planted
answer was built for Magic or for Netgen}, and an earlier draft of this
paragraph sat under a three-deck table and read as though one had been.
What stands in its place is weaker and is named as weaker: for Magic,
that the same invocation returned \num{39969214} boxes on the macros, so
the deck was reading geometry rather than skipping it; for Netgen, that
four of the seven sibling arms of the same comparison on the same design
return \emph{``Netlists do not match''}. Those are runs that came back red, not
controls with a known answer, and the difference is the one this paper
spends its length on.

\paragraph{The exclusions, which are not small.}
Metal density and antenna are \emph{separate top-level runsets} in this
PDK and neither was invoked, by this check or by the flow. Density bears
directly on the routing and fill this project drew, so ``zero KLayout
errors'' must not be read as including it. And the LVS result excludes
the macro internals by construction: what survives is the pin
\emph{list} (355 names on each of the four \texttt{2048x64} macros and
193 on each of the two \texttt{1024x32}, checked against the vendor's own
circuit netlist and identical as sets), the instance count, and the net
on all \num{1806} macro pin connections---a swapped data bit or a supply
strapped to the wrong rail would have shown, half a million bitcells
would not. An earlier draft printed \num{2130} for that last figure,
which is $6\times355$: the wide macro's pin list charged to all six
instances, when two of the six are the narrow one.

\paragraph{And the standard integrator's move is the wrong answer here,
which we did not expect.}
Section~\ref{sec:inventory} argued that abstracting a hard macro changes
\emph{what is looked at} rather than \emph{what counts as an error}. For
this deck and this macro that is \textbf{false}, and the measurement is
what says so. The abstracted run returns 123 error boxes where the full
run returns none outside the macros. \textbf{Three of the four rules
involved, covering 122 of the 123 boxes}---\texttt{M3.f} with 99,
\texttt{M4.f} with 16 and \texttt{M2.f} with 7---require metal wider
than \SI{10}{\micro\meter}, and
sizing the vendor macro's own geometry by $\pm\SI{5}{\micro\meter}$
leaves \emph{zero} surviving polygons on any of those layers---as does
the same test on this design's mask data in the exact windows where the
boxes were reported. The wide metal exists only in the abstraction,
because the macro's boundary description carries blanket obstruction
rectangles and the tool paints an obstruction as drawn metal.

\textbf{The remaining box is not covered by that argument, and an
earlier draft of this paragraph said ``every rule involved'' and swept
it in.} It is a single \texttt{M4.e} marker---\emph{Metal4 $>$
\SI{0.39}{\micro\meter} with runlength $>$ \SI{1.0}{\micro\meter}
spacing to unrelated m4}, a minimum-spacing rule whose width threshold
is \SI{0.39}{\micro\meter} and not \SI{10}{\micro\meter}---on a
\SI{0.02}{\micro\meter} by \SI{5.48}{\micro\meter} Metal4 sliver whose
lower-left corner the report gives, in micrometres, as
$(904.700, 1386.255)$, outside every macro's drawn
extent. What is known about it is that the full run over the real
geometry reports nothing at that coordinate, so it too is present only
in the abstracted view. What is \emph{not} known is a mechanism for it:
the sizing test that explains the other 122 is a test about wide metal
and says nothing about a rule that fires at
\SI{0.39}{\micro\meter}. One box in 123 does not change what the
comparison decides, and rounding it into the explained ones would have
been the move this section exists to object to.

So the seventeen-hour run reading real geometry is not the expensive
route to the same answer. It is the only one of the two that is right,
and an integrator who took the cheap route here would have shipped 123
violations that do not exist. We recommend nothing from this beyond the
measurement: it is one deck, one macro, one tool.

\subsection{The timing scope, which is the same defect inverted}

The system-on-chip fails 3{,}088 setup endpoints at the slow corner with
a worst slack of $-$\SI{4.6799}{\nano\second} and total negative slack of
$-$\SI{8267.18}{\nano\second}. It \emph{meets} setup at typical by
+\SI{4.3153}{\nano\second} and at fast by +\SI{9.5579}{\nano\second}, and
meets hold at all three. Alongside the setup failures the slow corner
carries 33 maximum-capacitance and 67 maximum-slew violations.

A nine-nanosecond spread between the corner that is met and the corner
that is not is the shape of a design optimised at one corner and checked
at another, and the configuration says that is what happened:

\begin{center}
\small
\begin{tabular}{lll}
\toprule
Key & Value & Binds \\
\midrule
\texttt{SETUP\_VIOLATION\_CORNERS} & \texttt{['*']} & the \emph{checker}
to all three corners \\
\texttt{MAX\_SLEW\_VIOLATION\_CORNERS} & \texttt{['*']} & the checker \\
\texttt{MAX\_CAP\_VIOLATION\_CORNERS} & \texttt{['*']} & the checker \\
\texttt{TIMING\_VIOLATION\_CORNERS} & \texttt{['*typ*']} &
\textbf{the repair} \\
\texttt{RSZ\_CORNERS} & \texttt{None} & \textbf{the repair} \\
\texttt{PL\_TIMING\_DRIVEN} & \texttt{False} &
\textbf{placement, entirely} \\
\bottomrule
\end{tabular}
\end{center}

The project had already found, at cost, that this flow's checkers examine
less than a reader assumes: its configuration file carries four separate
comment blocks documenting checkers whose shipped corner bindings made
them incapable of failing, including one where two checkers resolve an
empty corner list, warn, and let the flow exit zero. Those were found and
bound to \texttt{['*']}.

\textbf{The repair side was left at the default, and a default is not a
decision.} The checkers were widened to every corner; the optimiser that
is supposed to satisfy them was not. So the design is repaired against
typical, measured against slow, and the gap between them is reported as a
failure of the design rather than of the scope it was optimised at.

This is the same shape as Section~\ref{sec:placement} with the roles
exchanged. There, a check ran at a level above where the property could
fail. Here, a \emph{repair} runs at a corner below where the property is
checked. In both cases the artefact that would have shown it---a
placement, a slow-corner slack---exists on disk and was not consulted.

\paragraph{What the path itself says.}
The worst path launches at the processor's register-file read address and
ends at an SRAM macro input, passing through address decode and a long
exclusive-or chain that is the memory's error-correcting encoder. Its
delay is not evenly distributed. One \texttt{nand3} of minimum drive
strength on that path has an output transition of
\SI{1.611}{\nano\second} and contributes \SI{1.396}{\nano\second} of
delay while driving \SI{0.119}{\pico\farad}; a \texttt{nand2}, also
minimum strength, transitions in \SI{1.042}{\nano\second}. Two cells,
roughly \SI{2.2}{\nano\second}. \textbf{A minimum-strength cell driving
ten times its comfortable load is a resizer that did not run at that
corner}, and the 33 capacitance and 67 transition violations at the same
corner are the same fact counted a different way.

\paragraph{The scope hypothesis was wrong, and the flow's own source
says so.}
The paragraph above reads as though the corner split were the cause. It
is not, and we report that because the reasoning is the paper's subject.
\file{RSZ\_CORNERS} is unset, and the resizer step resolves
\texttt{RSZ\_CORNERS or STA\_CORNERS}---\emph{not}
\texttt{PNR\_CORNERS}---so it falls back to all three corners and had
been optimising the slow one all along. The resizer's own log confirms
it, loading the slow standard-cell and macro libraries.
\texttt{TIMING\_VIOLATION\_CORNERS} never reaches it: it feeds the
checker steps, where the explicit setting already overrides it.
\textbf{There is no one-key corner fix, and setting the key would have
narrowed what is optimised rather than widening it.} The table stays
because the scope split is real and worth a reader's attention; the
causal reading of it does not survive.

\paragraph{It is a plateau, not a cliff.}
Mean slack $-\SI{2.68}{\nano\second}$, median $-\SI{2.58}{\nano\second}$;
eight endpoints lie inside $-\SI{0.25}{\nano\second}$ and 87\,\% are
worse than $-\SI{1.5}{\nano\second}$. Repairing the worst handful buys
nothing. An independent static-timing sweep on the final netlist and
extracted parasitics reproduces sign-off exactly and closes the design
at \textbf{\SI{25}{\nano\second}}: $+\SI{0.3201}{\nano\second}$, zero
total negative slack, zero violating endpoints. The design's own clock
report had said as much---\texttt{period\_min = 24.68}.

\paragraph{Two hard floors, neither of them the optimiser's fault.}
The SRAM read is \SIrange{8.88}{9.29}{\nano\second} at the slow corner
from the vendor's own timing library, so \textbf{46\,\% of a
\SI{20}{\nano\second} budget is spent inside the macro before a single
gate}. And the cell library runs out: of the driver stages on the worst
thousand violating paths, \textbf{39\,\% are already at the largest
drive strength this library offers}, and design-wide 36.5\,\% of
standard cells belong to families with \emph{no} larger variant at all.
Only buffers and inverters reach the higher strengths.

Logic depth is not the wall. Substituting, on the worst path, the
fastest instance of each cell type seen anywhere in the report gives
\SI{8.801}{\nano\second} for its 73 stages against a
\SI{19.874}{\nano\second} budget---\textbf{the design runs about
$2.7\times$ slower than its own gates can go}, and what it is paying for
is load and the slew that load creates.

\paragraph{The violation counts were a red herring, and the inverse is
the finding.}
None of the 30 standard-cell slew violations and none of the five
standard-cell capacitance violations lies on any of the worst thousand
paths. Repairing all of them buys approximately nothing. \textbf{The
finding is that the limits are too loose to catch what hurts}: the only
transition limit in force is the library's own per-pin
\SI{2.5074}{\nano\second}, which is 12.5\,\% of the clock period in a
single transition, so the repair step legitimately left \num{4316} pins
alone. The gate we singled out earlier at \SI{2.4715}{\nano\second} sits
\emph{just under} that limit. The constraint was the problem, not the
tool.

\paragraph{The repair, run as a factorial, and it inverts its own
diagnosis.}
Two mechanisms were available: constrain transition and capacitance
explicitly, and make global placement timing-driven. They were run as a
$2\times2$ rather than bundled, because bundled they are unattributable.

\begin{center}
\small
\begin{tabular}{lrr}
\toprule
slow corner, \SI{20}{\nano\second} & placement not timing-driven &
timing-driven \\
\midrule
no explicit limits & $-4.6799$ (baseline) & $-3.3554$ \\
slew 0.75 / cap 0.10 & $\mathbf{-5.2145}$ & $-3.0306$ \\
\bottomrule
\end{tabular}
\end{center}

\textbf{The explicit limits alone made setup worse by
\SI{0.5346}{\nano\second}}, cost 6.98\,\% of standard-cell area and
1.29\,\% of routed wire, and pushed the minimum period from 24.68 to
\SI{25.21}{\nano\second}. \emph{All four arms, because printing two of
them was the thing this section is against:} minimum period 24.68 for
the baseline, 25.21 for the limits alone, \textbf{23.36} for
timing-driven placement alone and \textbf{23.03} for both, and the
period at which each closes with zero violating endpoints is
\SI{25}{\nano\second}, \SI{25}{\nano\second}, \textbf{\SI{24}{\nano\second}}
and \textbf{\SI{24}{\nano\second}}. An earlier draft printed the
baseline and the harmful arm and drew ``it did not close'' from them,
which is selective reporting inside the experiment offered as the
antidote to it. It was caught in audit, not in writing. The diagnosis that ranked them first had
modelled them at $+4.3$ to $+\SI{5.5}{\nano\second}$: \textbf{wrong by
about five nanoseconds and in the wrong direction}. Timing-driven
placement alone bought $+\SI{1.3245}{\nano\second}$; together they buy
$+\SI{1.6493}{\nano\second}$ on worst slack---but on total negative
slack, endpoint count, area, wirelength and power, \emph{timing-driven
placement alone is the better layout}, and its endpoint count is 2{,}102
against 2{,}746.

\textbf{It did not close at \SI{20}{\nano\second}}:
$-\SI{3.0306}{\nano\second}$ with 2{,}746 violating endpoints. It does
close at \SI{24}{\nano\second}, at $+\SI{0.9694}{\nano\second}$ with
none, and at \SI{23}{\nano\second} it is one endpoint short at
$-\SI{0.0306}{\nano\second}$. So the repair moved the period at which
this layout closes from 25 to \SI{24}{\nano\second} and did not reach
the target. What closes this design is the clock period, and
Section~\ref{sec:inventory} says what that means for the frequency this
paper does and does not print.

This is the paper's own method turned on the paper's own repair, and it
went against the diagnosis. A careful, measured analysis mis-ranked its
two candidate repairs by five nanoseconds and put the harmful one first.
The factorial caught it; a single bundled run would have reported
$+\SI{1.6493}{\nano\second}$ and credited it to both.

\section{Where the replicas actually are}
\label{sec:placement}

\subsection{The question nothing had asked}

The census of Section~\ref{sec:evidence} establishes that the mapped
netlist contains three banks of 55 flip-flops. It runs on the synthesis
netlist. Placement does not exist yet when it runs, so it cannot see
placement, and no other check in the project looked.

A majority vote over three registers placed in adjacent rows is defeated
by whatever corrupts two adjacent registers. Nothing in this flow asks
the placer to keep replicas apart, and nothing was asked.

\subsection{The instrument}

\file{hw/openlane/replica\_placement.py} reads the run's own final
design-exchange file and hierarchical netlist, attributes every mapped
flip-flop to a replica by instance path, attributes each to the storage
bit it holds where that net survives, and computes distances between cell
bounding boxes.

Its first assertion is that it counts the same flip-flops the synthesis
census counts. If the placed database and the netlist have come apart,
nothing else means anything. It reads 55, 55, 55---the census's own
figure on the same run.

Geometry, from the run's own files and the process design kit: die
\SI{1289.28}{\micro\meter} $\times$ \SI{313.74}{\micro\meter}; site
\SI{0.48}{\micro\meter}; row \SI{3.78}{\micro\meter}; the flip-flop is
\SI{12.96}{\micro\meter} $\times$ \SI{3.78}{\micro\meter}, which is 27
sites wide and exactly one row high.

\subsection{The measurement}

\begin{table}[htbp]
\centering
\caption{Replica placement in the frozen sign-off layout, 33{,}564
components. Every row is produced by one invocation of
\file{hw/openlane/replica\_placement.py}; six are additionally pinned in
\file{paper/claims.yaml}, and the per-replica count is registered for
replica~A alone, the tool reporting 55 for all three on the same run.}
\label{tab:placement}
\small
\begin{tabular}{lr}
\toprule
Flip-flops per replica & 55 / 55 / 55 \\
Cross-replica pairs & \num{9075} \\
\midrule
\multicolumn{2}{l}{\emph{The measurement of the mechanism}}\\
Centre-to-centre median, all cross-replica pairs & \SI{113.47}{\micro\meter} \\
\textbf{Centre-to-centre median, corresponding bits} &
\textbf{\SI{27.52}{\micro\meter} --- 4.1$\times$ closer} \\
\midrule
\multicolumn{2}{l}{\emph{Absence of segregation, against its own null model}}\\
Nearest fellow flip-flop is in another replica & 103 of 165 \\
\quad 200 label permutations of the same positions & $110.73 \pm 8.26$ \\
\quad analytic expectation & $110.67$ \\
\quad so the observed value is & \textbf{0.94\,sd \emph{below} chance} \\
\midrule
\multicolumn{2}{l}{\emph{Distances, which are floors as much as findings}}\\
Centre-to-centre, minimum & \SI{3.78}{\micro\meter} --- one row, and the
floor: the cell is \SI{12.96}{} by \SI{3.78}{\micro\meter} on a
\SI{3.78}{\micro\meter} row pitch \\
Edge-to-edge, minimum & \SI{0.00}{\micro\meter}, structurally guaranteed
for any two cells in adjacent rows whose spans overlap \\
Pairs whose cell outlines touch & 49 (41 in adjacent rows, 8 in the same row) \\
Pairs within one \SI{0.48}{\micro\meter} site & 51 \\
Corresponding-bit pairs, minimum & \SI{3.78}{\micro\meter} \\
Corresponding-bit pairs touching & 18 (a lower bound; see below) \\
\bottomrule
\end{tabular}
\end{table}

The closest pairs are the same bit of the same word. Two replicas of
configuration bit 35 sit in adjacent rows at the identical horizontal
span, sharing all 27 sites of width; bit 32 is the same case.

\textbf{The 103-of-165 row needs a baseline, and with one it is not the
finding it looks like.} Shuffling only the replica \emph{labels} over the
same 165 placed positions, keeping the 55/55/55 partition, gives
$110.73 \pm 8.26$ over 200 permutations against an analytic expectation
of $110.67$: the observed value is \textbf{0.94 standard deviations
\emph{below} chance}. So the statistic does not say the placer drew the
replicas together. It says the three banks are interleaved to the point
of being indistinguishable from an arbitrary assignment of labels to
these positions---which is a finding, and a different one. Three
separated banks would give a value near zero, many standard deviations
away, and that is what it would detect.

\textbf{The row that measures the mechanism is a ratio.} A shared voter
ties the three copies of one bit together and nothing else, so if
wirelength minimisation is what co-locates them the effect belongs on
\emph{corresponding} bits rather than on cross-replica pairs at large.
It is there: \SI{27.52}{\micro\meter} against
\SI{113.47}{\micro\meter}, \textbf{4.1 times closer}. The instrument had
printed both on every run since it was written; an earlier draft
asserted the mechanism in bold and omitted the number in its own tool's
output that evidences it.

\subsection{The mechanism, which matters more than the numbers}

This is not a bug and not an unlucky seed. \textbf{The placer is
minimising wirelength to a shared voter, and that objective pulls the
three copies of each bit together.} It is what was asked of it.

One asymmetry was not designed and is worth recording with its size.
Per pairing the minima are A--B \SI{3.78}{\micro\meter}, A--C
\SI{4.75}{\micro\meter}, B--C \SI{3.90}{\micro\meter}, and of the 49
abutting pairs 31 are A--B, 11 A--C and 7 B--C; the medians are
indistinguishable, so the whole effect sits in the abutment tail. C
carries the mixing transform of Section~\ref{sec:mechanisms}---and C is
also the most internally clustered bank, and no build with that
transform disabled was ever placed. \textbf{We record the asymmetry and
do not attribute it}, for the same reason
Section~\ref{sec:determinism} refuses to credit a clock-gating
improvement to a repair: a rule suspended when a result goes your way is
not a rule. In any case it buys nothing that matters---A--B, the pairing
the transform does not touch, is the worst of the three, and a vote
needs only two replicas corrupted.

\subsection{What the tool undercounts, and says}

Fifteen of replica C's flip-flops cannot be attributed to a bit index,
because C stores a transform of the value and the corresponding nets may
be optimised away. The corresponding-bit figures are therefore a
\textbf{lower bound}, and the tool prints the word \textsc{undercount}
when it reports them. A deeper one-off analysis that resolved the
remaining fifteen through the decode network found 165 corresponding-bit
pairs, 28 of them touching, and two of three voter inputs in contact on
30 of the 55 bits. \emph{That analysis was run once by hand, is not
committed, is reproducible from no script in this tree and is asserted
by no test, and the three figures are printed to bound the direction of
the undercount and for nothing else.} The committed tool reports the
conservative figure and prints the word \textsc{undercount}, because a
lower bound quoted as a measurement is the failure this paper is
about---and printing those three unmarked, one sentence before saying
so, was that failure.

\textbf{An earlier draft added that those three were ``the only figures
in this section a reader cannot re-derive'', and that sentence does not
survive being checked against what the tool prints.} Read against the
output of \file{hw/openlane/replica\_placement.py} and against
\file{paper/claims.yaml}, the two between them do cover the per-replica
counts, the cross-replica pair count, both medians, the minimum
centre-to-centre and edge-to-edge distances, the abutting and
within-one-site counts, the corresponding-bit minimum and abutment, the
103-of-165 statistic with both its permutation and its analytic null,
the fifteen unattributed flip-flops, and the die, site, row and cell
geometry. What they do not cover, and no other file in this repository
covers either, is the per-pairing minimum A--C \SI{4.75}{\micro\meter}
and the 31/11/7 split of the 49 abutting pairs across the three
pairings---the tool breaks out neither, its list of closest pairs
happens to expose the A--B and B--C minima, and there is no per-pairing
table anywhere---together with the statement that C is the most
internally clustered bank, which no committed figure supports. The three
system-on-chip per-bit separations in the next subsection sit one tier
down rather than outside: they are written into the project's record,
but the attribution method that produced them is not committed, as that
subsection says, so a reader can read them and cannot re-derive them.
The figures in the last
subsection---core utilisation, unoccupied area, the longest routed net,
the pilot's slack and the cost of the one placement constraint this
project has ever applied---are read from run reports and from documents
cited in the text, which is the tier
Table~\ref{tab:inventory}'s caption describes and not this section's
instrument. \textbf{So the count was wrong and the shape of the error is
this paper's own: a sentence bounding what a reader cannot check,
written without checking it against the tool it was bounding.}

\subsection{It is not one layout}

The same picture appears on the system-on-chip across three independent
triplicated structures---a 29-bit watchdog bank, a 23-bit boot bank, a
25-bit accelerator bank---on two independent hardening runs, by a
completely different attribution method, because that flow flattens
instance names and the cells must be reached through fan-in cones.
Per-bit critical separations are \SI{4.48}{\micro\meter},
\SI{7.58}{\micro\meter} and \SI{3.90}{\micro\meter} minimum, with
different-replica cells holding the same output bit abutting in the same
row. Two designs, three layouts, three attribution methods---only one of
which is committed---and one answer.

\subsection{What this does not say, in the paper's own voice}

\textbf{A distance is not a safety argument, in either direction.} No
single-event or multiple-bit-upset cross-section, charge-collection
radius or linear-energy-transfer threshold for this process was measured
and none is available: the PDK publishes no radiation data at all. These
are cell bounding boxes, not sensitive nodes---the storage node inside a
\SI{12.96}{\micro\meter} $\times$ \SI{3.78}{\micro\meter} rectangle is
somewhere in that rectangle and this work does not say where.

So the measurement establishes that this layout \emph{permits} a
common-mode upset of two replicas across a large fraction of the
protected word. It does \textbf{not} say one occurs at any fluence, and
it does \textbf{not} say that some other separation would prevent one.
The sentence ``the replicas are a few micrometres apart, so a single ion
track defeats the vote'' does not appear in this paper, because it
silently supplies a conversion from micrometres to a coincidence
probability that we do not have.

What it does establish, and what it obliges, is a change to the wording
of every claim this project makes about the state in question.
\textbf{``Protected by triple modular redundancy'' is unsupported at the
die and must read ``triplicated \emph{in the netlist}''.} Note the verb:
the measurement does not make those sentences \emph{false}---that would
need exactly the conversion from micrometres to a coincidence probability
that the paragraph above forbids---it removes their evidence. That is a
correctness change and not a style preference.

An earlier draft of this paragraph said the corpus ``now reads
`triplicated in the netlist' throughout''. It did not: the sweep was
announced in the sentence that was supposed to record it, which is this
paper's own subject committed in the sentence claiming the fix. The
sweep has since been made, across the project's front page, both
datasheets, four hardening documents, two register-transfer headers and
two test docstrings; \file{paper/claims.yaml} carries an entry that
greps for the unqualified phrase, so the claim is now checked rather
than asserted.

\subsection{A constraint exists, and it is not free}

Three mechanisms could impose separation, and this project has exercised
two of them before: a per-instance manual placement hint, which the
pilot flow already runs as a step and which does nothing when unset; a
region-and-group construct in the placer's database, which fences a group
into a rectangle; and soft placement obstructions, which can evacuate a
rectangle but cannot fence a group into one and are useless alone.

The price is known from the one time this project constrained placement,
on an unrelated block: the region form \emph{did not route}, and the
manual-placement form routed but cost roughly \SI{7}{\nano\second} of
worst negative slack at the intermediate stage where it was first
measured---\SI{4.88}{\nano\second} at sign-off---and moved the problem
rather than solving it,
because a constraint that captures one half of a co-placed cluster moves
the crossing to the cut. For triplication that means each replica must
move together with its own fan-in and fan-out cone, not merely its 55
flip-flops.

On the pilot the conditions are friendlier---no macros, 48.75\,\% core
utilisation, \SI{201400}{\micro\meter\squared} of the core unoccupied
against \SI{2694}{\micro\meter\squared} for one bank, and a design that
already routes a \SI{901.08}{\micro\meter} net and closes with
\SI{1.2262}{\nano\second} to spare. The binding constraint is timing, not
area. We did not attempt it: the design is frozen for a shuttle, and an
unmeasured constraint applied to a frozen design is not an improvement.

\section{When the evidence was wrong}
\label{sec:evidence}

\subsection{The claim}

A fault-tolerance property is not one claim. It is a claim at the source,
a claim in the mapped netlist, and a claim in the die, and the three can
disagree. The instrument that establishes it at one level is
\emph{structurally incapable} of refuting it at the next---not merely
unlucky, not merely incomplete, but built at a level where the failure
cannot be expressed.

Table~\ref{tab:evidence} is the paper's contribution. Two rows, both from
this project's own published claims, both reproducible from committed
artefacts, both carrying a mutation that has been observed to make the
replacement instrument fail.

\begin{table}[htbp]
\centering
\caption{Two levels at which a redundancy claim was false while a check
aimed at it passed.}
\label{tab:evidence}
\small
\begin{tabularx}{\textwidth}{L{2.4cm}XX}
\toprule
& \textbf{R1 --- collapsed in the netlist} & \textbf{R2 --- reset by its
own correction} \\
\midrule
\textbf{Defect} &
Three declared replicas became one physical bank; the voter read it three
times &
The voted word decoded combinationally onto the asynchronous reset of
every flip-flop \\
\addlinespace
\textbf{Evidence the property was absent} &
At 4$\times$2: 362 references to \texttt{cfg\_a}, \textbf{0} to
\texttt{cfg\_b}, \textbf{0} to \texttt{cfg\_c}; \num{1045} mapped
flip-flops against \num{1161} declared. Reproduced on a second PDK &
Reset-cone census: of five asynchronous-reset fan-in source sets, one has
a combinational fan-in of \textbf{25 flip-flops, all replica bits}.
Repaired netlist: 26 sources, the extra one being the register now in
between \\
\addlinespace
\textbf{Blind instrument 1} &
The register-transfer fault-injection campaign aimed at this TMR:
5 injections per replica, all \textsc{corrected}, 0 \textsc{sdc} ---
\textbf{bit-identical before and after the repair} &
The gate-level classifier: 174/174 \textsc{corrected}, no announcement,
\textbf{no reset column at all}. It samples on the clock edge and cannot
see a sub-clock pulse \\
\addlinespace
\textbf{Blind instrument 2} &
A flip-flop-total A/B with and without a preservation attribute:
\textbf{the same total both ways}, read as ``nothing merged''. It is
equally consistent with ``the attribute does not work'', which was true &
$k$-induction on
\texttt{rst\_req\_o = in\_reset}: \textbf{true of the defective and the
repaired design alike}, because they are the same function of the same
state at that level \\
\addlinespace
\textbf{Catching instrument} &
Per-replica census of mapped flip-flop cells \emph{by instance path}, in
two independent technology mappers, with every preservation attribute
deleted from the source text, re-run on the sign-off netlist: 55/55/55 &
An \emph{event} counter on the reset pin; a satisfiability query over the
mapped cells; a whole-netlist census of asynchronous-reset fan-in \\
\addlinespace
\textbf{Detector proved able to fail} &
Pre-repair source through today's guard: \textbf{4 failed}, census
$\{0,0,0\}$ at \num{1045} flip-flops in both mappers. A
\emph{functionally identical} mutation (removing one replica's transform)
fails 3 of 40 &
One byte-identical 174-injection plan on four netlists:
$1.0000 \to 0.0000$ resets per correction, every masking field identical.
Two committed mutations both fail the proof \\
\addlinespace
\textbf{Cost of the repair} &
\num{1045} $\to$ \num{1155} flip-flops;
\SI{106174}{\micro\meter\squared} $\to$
\SI{109058}{\micro\meter\squared} (+2.7\,\%) &
\textbf{One flip-flop.}
\SI{643.0536}{\micro\meter\squared}, +0.0993\,\%; die area unchanged \\
\bottomrule
\end{tabularx}
\end{table}

\subsection{Two rows, not six, and why the difference is the method}

We began with nine candidates and tested each against four conditions:
the defect was real; \emph{something was checking and passed}; a distinct
detector exists in the tree; and that detector has been observed to fail
on a defective input. Seven did not survive, and the reasons are more
useful than the survivors.

\paragraph{Absent is not blind.}
Four candidates failed because nothing had been checking. A windowed
watchdog was recommended, built, measured at zero catches in thirty and
removed at source---a good decision, correctly taken, and \emph{not} an
evidence failure, because the mechanism did not exist when it was
recommended. A shipped-netlist census caught a real gap the first time it
ran, which is coverage arriving, not evidence failing. Conflating the two
would be this paper's own thesis committed in its own text: ``no one
looked'' and ``someone looked and could not see'' are different claims
and only the second is interesting.

\paragraph{A red result read too wide is still a red result.}
One candidate was a timing violation attributed to the wrong logic cone.
The analysis had reported both the violation \emph{and} the correct
launching register; the error was in reading it. That is a real mistake,
recorded and corrected in the project, and it is not an instrument being
blind.

\paragraph{Tightening a check is not a new instrument.}
Several candidates were the same check with a narrower window. They
belong in a changelog.

\paragraph{And one candidate was no defect at all.}
A census panic that appeared to show undercounted replicas turned out to
be the panic's own instrument counting by net name instead of instance
path. \textbf{No published replica count was ever wrong.} We report this
because a paper cataloguing false alarms that itself raised one, and
suppressed the fact, would be worthless.

\subsection{The third level, found by turning the method on itself}

Section~\ref{sec:placement} is not in Table~\ref{tab:evidence}, and it
is the strongest thing in this paper.

The catching instrument for R1 is a census of mapped flip-flops. It
passes. It is mutation-checked, runs in two technology mappers and
survives the deletion of every preservation attribute. And it establishes
\emph{logical} redundancy only, because it runs before placement exists.
The claim it supports was being read as a claim about the manufactured
part, and at that level it is false.

We did not find this by being careful. We found it because a hostile
reading of the project's own record asked which measurement, if a
reviewer made it first, would sink the argument---and then we made it.
The answer came back badly, and it is published because a paper arguing
that green checks do not establish what they appear to establish is
\emph{confirmed}, not damaged, by discovering that its own green check
did not.

\subsection{The middle link, which did not exist}

A census counts cells and walks cones. It never establishes that the
netlist \emph{computes} what the source says. A synthesis fault that
preserved three banks of 55 flip-flops and miscomputed the voter would
pass every check in this project---and the voter is the one module a
flip-flop census structurally cannot check, because it has no flip-flops.

Until this paper there were \textbf{zero} equivalence checks in the
repository, while the tool sat in the pinned toolchain the whole time.
There are now three, on the modules chosen for that reason: the voter at
its instantiated width, and both SEC-DED codecs. The gate side is
\emph{extracted from the sign-off run}, not synthesised fresh, so the
result is about the artefact that is frozen rather than about today's
recipe. Two negative controls---one term dropped from the majority
function, one bit flipped in a parity row---are part of the default
target, and a control that \emph{passes} is treated as an error.

The equivalence checker itself produced a finding. It partitions at
module outputs and feeds each partition's cut points as independent free
inputs, which for the decoder manufactures an unreachable
state---a syndrome no codeword produces---and reports a counterexample
\textbf{on a correct netlist}. Binding the cut points to the reference
design's own nets removes the freedom without weakening the proof. The
general statement is worth carrying: \textbf{partitioned equivalence
checking is sound for proofs and incomplete for refutations}, and a gate
written as \texttt{eqy \ldots{} \&\& echo ok} would have failed on a
correct netlist on the first day it ran.

\subsection{The generalisable form}

\begin{quote}
For each level of representation between the source and the die, a
fault-tolerance property is a separate claim. Establish it at each level
with an instrument that lives at that level. An instrument that lives
above the level at which the property can fail cannot refute it, and will
return green. Every replacement instrument must carry a mutation that has
been observed to make it fail; a guard that has never returned red is not
a guard.
\end{quote}

Concretely, for triplication in a standard-cell flow: check the source by
reading it; check the netlist by censusing mapped cells \emph{by instance
path}, in more than one mapper, with preservation attributes deleted;
check the \emph{function} by equivalence checking, because a census
cannot; and check the die by measuring where the cells are. We did the
first three. We did the fourth for the first time while writing this
paper, and it failed.

\section{Flow determinism, and what it licenses}
\label{sec:determinism}

Several cost figures here are differences between two runs of the same
flow. That difference means nothing unless the flow is deterministic, so
we measured it---on \emph{one} netlist, the system-on-chip's protected
boot block, through \emph{one} 62-step flow configuration on one host.
It licenses the cell-count deltas of Section~\ref{sec:mechanisms} and
the four-arm comparison of Section~\ref{sec:w9}. It is \emph{not} a
measurement of the pilot's sign-off flow, which was never re-run, so the
pilot figures in Table~\ref{tab:inventory} and the layout of
Section~\ref{sec:placement} carry no repeatability evidence of their
own. An earlier draft's section title generalised one netlist through
one configuration into ``every delta in this paper''.

The same netlist through the same 62-step flow twice: the sign-off
metrics file is byte-identical on all 199 keys; the routed design
exchange file, the database, both netlists, the constraints, three
timing libraries, the abstract layout and the extracted netlist are
byte-identical; the mask data, the layout database, the parasitics and
three delay files are identical \emph{modulo the date stamps their
writers embed}; and 58 of 62 steps are identical in every metric and
every view.

Under the vocabulary of Lamb and Zacchiroli~\cite{lamb2022} this is
\textbf{repeatability}, not reproducibility: one host, one operating
system, one tool build, and not bit-identical because of embedded
timestamps. We say so in those words.

Its value is instrumental rather than headline. It converts every
``noise floor'' this project had previously reported into
\emph{netlist sensitivity}, and it is what licenses the cell-count deltas
of Section~\ref{sec:mechanisms} and the four-arm comparison of
Section~\ref{sec:w9} to be read as measurements at all. Without it, a
$1.0000 \to 0.0000$ difference across two designs would be a difference
between two runs.

It also has teeth in the other direction, and the project has been cut by
them. Two clock-gating checks improved after a repair, and the
improvement was \emph{not} claimed for the repair, because the module
that carried one of them was byte-identical between the two builds: its
check closed because the netlist landed differently. A delta between two
layouts of this flow is a delta of netlists, and that rule cost the
project a headline it would otherwise have printed.

\section{Limitations and threats to validity}
\label{sec:limits}

Most of these are in the sections they belong to. They are collected here
because a reader is entitled to find them in one place.

\paragraph{The evidence tier.}
No silicon, no beam, no total-ionising-dose curve, no single-event
cross-section, no latch-up data, and no radiation characterisation of the
process at all. Under the space product-assurance standard this is not an
evidence tier that qualifies anything~\cite{ecss60_15c}. Work that does
reach it exists for spiking accelerators~\cite{nijsink2026}, and this
paper is explicitly below it. What gate-level fault injection can bound
is logical and architectural masking and the structural divergence
between representation levels; what it cannot touch is multiple-bit
upset, transient width, latch-up, dose and spatial correlation.

\paragraph{No rate, anywhere.}
Every proportion is conditional on an upset having landed in the
injection window, uniform over targets and cycles. Simulated time is not
exposure.

\paragraph{Sample size.}
Proportions are largely reported without confidence intervals derived
from statistical fault-injection theory~\cite{leveugle2009}. This is a
correctable weakness and it is not corrected here.

\paragraph{Fault model.}
Single-bit, flip-flop only, one workload per campaign. No transients, no
multiple-bit upsets, no permanent faults. Combinational cells are not
injection sites, so a 528-cell decoder is outside every campaign in this
paper.

\paragraph{The proofs are scoped to one fault.}
Every proof about a voted structure is scoped to one bad replica at a
time. No proof here can express the failure mode that
Section~\ref{sec:placement}'s geometry makes relevant: two of three
replicas corrupted.

\paragraph{Physical independence is measured and absent.}
Section~\ref{sec:placement}. And the measurement itself is of cell
bounding boxes on one layout per design, with no cross-section to
interpret it against.

\paragraph{Equivalence is three modules, not the design.}
Whole-design equivalence against the placed netlist was attempted and
reached roughly half the partitions; no figure for it is printed here,
because ``no failing partition was traced to a synthesis defect'' is
strictly weaker than ``the netlist is correct'' and we decline to print
the weaker statement as though it were the stronger one.

\paragraph{The design under test is not the specified design.}
Section~\ref{sec:design}. Every figure carries its geometry.

\paragraph{No application result.}
No accuracy on any task, no benchmark, no compiler, and one small
workload. One accuracy figure that this project once carried has been
withdrawn: it was unsupported in both of the readings available for it,
and re-checking the sibling design's own repository found the number
occurs nowhere in it.

\paragraph{The verdicts are not in the repository.}
Section~\ref{sec:formal}. A reader who clones obtains the jobs and the
harnesses, and must re-run to obtain a single verdict. The run trees that
back the fault-injection campaigns are likewise build output; digests are
published in their place, and a digest establishes identity, not
correctness.

\paragraph{The continuous integration, and a correction we made to
ourselves twice.}
Every result in this paper was obtained locally. What continuous
integration adds is independent re-execution, and this project has not
had it reliably: of 61 workflow runs since 3 September 2026, 57 failed.
Eight, all from 10 September, were \emph{refused before starting} on a
billing annotation. The rest executed and failed for real reasons, and
one of them was a missing Python package in the workflow itself, which
made the licence check fail eight times without ever reaching the check.

We report the sequence because it is this paper's subject twice over.
First: three documents and this paper's own source header said those
checks ``run in CI'', which described a configuration and not an
observed event. Then, correcting that, we wrote that \emph{every} run
since 3 September had been refused before starting---taking one day's
annotation and applying it backwards across a week. That is a red result
read wider than what it looked at, which is the failure this paper
catalogues, committed in the sentence that announced finding it
elsewhere. The checks now live in one script that does not need a
runner, and the workflow calls it.

\paragraph{Who checked the checker.}
Every audit reported here was performed by the author of the thing
audited. Our answer is that an audit is worth something only if it is
mechanically re-runnable and mutation-checked, and Appendix
\ref{app:artefacts} is one command per claim. The project's own sharpest
statement against itself is worth more than any protestation of rigour:
\emph{a mutation that passes is a mutation that was testing the wrong
thing}---which is how a preservation attribute was found to preserve
storage that was already dead.

\paragraph{Four references we could not verify.}
A 2021 study of triplicated-layout single-event tolerance, an automated
triplication-flow paper, and two open fault-injection frameworks surfaced
in the literature search and could not be corroborated to author, venue
and year. They are not cited. We name them here because ``the literature
contains X'' is itself a claim, and an unverified one belongs with the
other unverified things.

\section{Related work}
\label{sec:related}

Most of what this design does is not new, and the honest way to open a
related-work section is by saying which parts.

\paragraph{Spiking accelerators on open PDKs with open flows are done.}
OpenSpike~\cite{openspike2023} taped one out in SkyWater 130\,nm at
roughly thirty times this design's scale in 2023, and further open
130\,nm neuromorphic silicon followed~\cite{kumaresan2026}. An SNN on an
open PDK is not a contribution and is not claimed as one. The
neuromorphic-for-space framing is a survey topic~\cite{naoukin2023}.

\paragraph{Neuromorphic hardware under radiation has beam data.}
Nijsink~\emph{et~al.}~\cite{nijsink2026} put an open SNN
processor~\cite{frenkel2019} in a neutron beam and extracted
cross-sections. That is the tier above this work, and the difference is
not one of degree.

\paragraph{Triplication and its synthesis survival are old.}
Voter placement, not triplication itself, dominates residual
cross-section in the FPGA literature~\cite{kastensmidt2005}; the
synchronisation-voter algorithms are standard~\cite{johnson2010};
post-synthesis insertion tools exist because logic optimisation removes
source-level redundancy~\cite{tmrg2017}; and vendor guidance documents
the preservation attributes~\cite{esa_ftmr}. Our contribution on this
axis is not that the merge happens. It is the measurement of what did not
notice, and the bound that says how wide a word must be before a third
transform exists.

\paragraph{Identical replication is known to be weak.}
Mitra~\emph{et~al.}~\cite{mitra2002} is the citation, and the diversity
line in silicon~\cite{safede2021,tcls2016,tcls2019} operationalises it.
Section~\ref{sec:placement} is a physical instance of the same idea:
identical replicas that the placer is free to co-locate are correlated by
geometry, whatever their logic says.

\paragraph{Hardened RISC-V systems with fault injection exist, with
silicon.}
Trikarenos~\cite{rogenmoser2023} combines ECC-protected memory with
triple-core lockstep and characterises it in a beam \emph{and} in
gate-level injection; on-demand redundancy grouping~\cite{odrg2022}
explores the reconfigurable version. This is the work our fault-injection
methodology is compared against, and the comparison is unfavourable on
evidence tier.

\paragraph{The closest published work to this paper's angle.}
Santos~\emph{et~al.}~\cite{santos2023} state their contribution as
extending fault \emph{observability} in a hardened RISC-V
system-on-chip---closing the gap between ``the system corrected
something'' and ``we can say what it corrected''.
Section~\ref{sec:w9} is exactly that gap and is one level deeper: the
system could say what it corrected and could not say what the correction
cost.

\paragraph{Cost optimality.}
CLEAR~\cite{clear2018} searched 586 cross-layer resilience combinations
against cost and residual silent-data-corruption rate. We produce no such
curve and therefore make no cost-optimality claim.

\paragraph{Codes, scrubbing, ECC proofs.}
Hsiao~\cite{hsiao1970} for the code; Saleh~\emph{et~al.}~\cite{saleh1990}
for scrub-interval reliability; Kumar~\cite{kumar2023} for the formal
proof of an ECC decoder, which is routine industrial practice and not a
contribution.

\paragraph{Hash-pinned artefacts.}
Lamb and Zacchiroli~\cite{lamb2022} is the framework. We found no
published application of it to a PDK, an open flow and a mask-data
freeze, and no prior work publishing a re-runnable post-place-and-route
triplication-survival check on an open flow. That, and the placement
measurement, are the parts of this paper we believe are not already in
the literature.

\section{Conclusion}
\label{sec:conclusion}

We set out to show that a design's configuration state was protected by
triple modular redundancy. At the source it was. In the mapped netlist it
was not, and the campaign built to detect exactly that produced identical
output on the broken design and the fixed one. We repaired it and built
an instrument that lives at the netlist. At the die we cannot show it is
protected: the replicas are not separated, nothing ever established that
they would be, and the instrument we built runs before the die exists and
could not have told us either way.

Separately, a protection mechanism that worked perfectly---174 of 174
upsets corrected, no error, no wrong answer---restarted the machine every
single time it worked, through a path the register-transfer level cannot
express and a classifier that samples on the clock edge cannot see. One
flip-flop, and a census of what feeds an asynchronous reset.

There is no silicon here and no beam, and this paper does not pretend to
the evidence tier that would need. What it offers instead is the missing
rung below: for each level between a source file and a die, an instrument
that lives at that level, and a mutation proving the instrument can fail.
The two places where ours could not are published above, and the third is
the one we found by asking which measurement would embarrass us most.

\section*{Reading the bibliography}

Each entry's note carries a verification key. \textbf{[V]} means
existence, authorship and venue were corroborated by an authoritative
source; \textbf{[P]} means existence and venue were corroborated and one
metadata field---a DOI, a page range, part of an author list---was
reconstructed from convention or taken from a secondary source and not
independently confirmed. References that reached only \emph{unverified}
are not cited here at all; four are named in
Section~\ref{sec:limits} instead, because ``the literature contains
$X$'' is itself a claim and an unverified one belongs with the others.

\bibliographystyle{plain}
\bibliography{refs}

\appendix

\section{Artefacts and how to re-derive every number}
\label{app:artefacts}

Everything below is in the public repository. Run trees and simulation
records are build output and are not tracked; digests are published in
their place, and a digest establishes identity, not correctness.

\begin{verbatim}
# the numbers this paper registers -- not every number it prints
python3 paper/check_claims.py --list
python3 paper/check_claims.py

# the placement measurement of "Where the replicas actually are"
python3 hw/openlane/replica_placement.py \
        hw/openlane/pilot_ihp/runs/signoff-6x2

# its guard, which pins the measured values rather than demanding
# separation, because the design is frozen
.venv/bin/python -m pytest sw/tests/test_replica_placement.py -q

# equivalence, with both negative controls; a control that PASSES
# is an error
make -C formal/eqy

# the netlist censuses of "When the evidence was wrong"
.venv/bin/python -m pytest sw/tests/test_synthesis_guards.py -q

# the suites, recorded rather than asserted
scripts/verify.sh
\end{verbatim}

\section{Claims withdrawn while this paper was written}
\label{app:withdrawn}

A ledger with no open rows is not a ledger. Each of these was a
statement this project had published and could not support.

\begin{itemize}\itemsep2pt
\item \textbf{A classification-accuracy figure} attributed to a sibling
design, tagged as a verified result. It collides with a placement
utilisation of the same value for the same design, and re-checking that
design's own repository shows the number occurs nowhere in it. Both
readings are unsupported. No accuracy figure is claimed anywhere in this
work.
\item \textbf{A compute-ratio comparison} against the reference product,
derived from a figure listed under a heading naming a two-page brief
while citing the vendor's product page. The figure is on the page and not
in the brief; the sourcing was wrong, not the figure. An independent
review had raised this and it was rejected on a re-verification that
named no artefact. The list is now split by source and the artefact is a
digest of the page as fetched.
\item \textbf{A block-area cost} for the one-flip-flop repair of
Section~\ref{sec:w9}. It does not reproduce across recipes and under the
project's own driver the sign inverts. Only the flip-flop count and the
whole-design area are printed.
\item \textbf{A busy-power figure} attributed to added clock gates,
withdrawn along with its own replacement, because across three netlists
the combinational group ranges too widely for the gates to order it.
\item \textbf{Two suite totals} recorded by overlapping runs of the
verification script, which counted each other's output: 558 and 624
cocotb tests for a suite of 401. The rows stand in the log with a
correction beside them, because a log that loses its own bad measurements
is not a log.
\item \textbf{A licence checker's own compliance.} It searched a
character window for its identifier string, and its own prose contains
that string, so it reported itself compliant while carrying no header.
Narrowed to a comment line in the first twelve.
\end{itemize}

\end{document}
